\begingroup
\renewcommand{\arraystretch}{1.45}
\clearpage\fancyhead[L]{\small LEMNISKATISCHE FUNKTIONEN II.}\pdfbookmark[0]{Über die Lemniskate}{de-curva-lemniscata}\SourceLeaf{422}

\begin{center}
\textbf{ÜBER DIE LEMNISKATE.}\par\medskip
\rule{2cm}{0.3pt}
\end{center}

\CurvaArticle{1}{\textbf{1.}}
Setzt man das Integral $\int\frac{\mathrm{d}x}{\surd(1-x^4)}$ von $x=0$ bis $x=s$, $=\varphi$, so nennen wir $s$ den lemniskatischen Sinus von $\varphi$, $s=\sin\operatorname{lemn}\varphi$.

\Needspace{34\baselineskip}
\CurvaArticle{2}{\textbf{2.}}
Der Wert des Integrals von $x=0$ bis $x=1$ ist nach \textsc{Stirling} $=1.3110287771\;4605987$; wir haben diesen Wert unter Verwendung der Formel $\operatorname{arc}\sin\operatorname{lemn}\frac{7}{23}+2\operatorname{arc}\sin\operatorname{lemn}\frac12$ bis zur elften Ziffer als richtig gefunden (\textsc{Euler} hat $1.311031$). Für die Potenzen dieser Zahl, deren Doppeltes wir stets mit $\varpi$ bezeichnen werden, haben wir Folgendes gefunden
\Needspace{24\baselineskip}
\[
\begin{array}{r@{\quad}c@{\quad}r@{\quad}l@{\quad}l}
1& \mathord{.}\mathord{.}\mathord{.}\mathord{.} &1.3110287771&4605990680&320.7\\[5pt]
2& \mathord{.}\mathord{.}\mathord{.}\mathord{.} &1.7187964545&0509311.7&\\[5pt]
4& \mathord{.}\mathord{.}\mathord{.}\mathord{.} &2.9542612520&1927863.4&\\[5pt]
5& \mathord{.}\mathord{.}\mathord{.}\mathord{.} &3.8731215170&0712625\;4&\\[5pt]
6& \mathord{.}\mathord{.}\mathord{.}\mathord{.} &5.0777737656&5251025.3&\\[5pt]
8& \mathord{.}\mathord{.}\mathord{.}\mathord{.} &8.7276595451&8251569.0&\\[5pt]
9& \mathord{.}\mathord{.}\mathord{.}\mathord{.} &11.4422128208&59&\\[5pt]
12& \mathord{.}\mathord{.}\mathord{.}\mathord{.} &25.7837864151&41749&\\[5pt]
13& \mathord{.}\mathord{.}\mathord{.}\mathord{.} &33.8032859402&5&
\end{array}
\]


\SourceLeaf{423}
\Needspace{12\baselineskip}
\[
\begin{aligned}
\operatorname{log}\operatorname{brigg}\dfrac12\varpi&=0.1176122226\quad9692.2\\[5pt]
\operatorname{log}\operatorname{hyp}\dfrac12\varpi&=0.2708121550\quad7159155410\quad6425\\[5pt]
\operatorname{log}\operatorname{hyp}\varpi&=0.9639593356\quad3153686352\quad36577
\end{aligned}
\]
Den lemniskatischen Sinus von $(\tfrac12\varpi-a)$ werden wir den lemniskatischen Kosinus von $a$ nennen.

\Needspace{22\baselineskip}
\CurvaArticle{3}{\textbf{3.}}
Für die Gleichung
\Needspace{8\baselineskip}
\[
\dfrac{\mathrm{d}x}{\surd(1-x^4)}+\dfrac{\mathrm{d}y}{\surd(1-y^4)}=0
\]
findet sich das folgende vollständige Integral
\Needspace{12\baselineskip}
\[
\dfrac{x\surd\dfrac{1-yy}{1+yy}+y\surd\dfrac{1-xx}{1+xx}}
{1-xy\surd\dfrac{1-xx}{1+xx}\,\dfrac{1-yy}{1+yy}}=C
\]

Ein Integral derselben Gleichung ist aber
\Needspace{8\baselineskip}
\[
\operatorname{arc}\sin\operatorname{lemn}x+\operatorname{arc}\sin\operatorname{lemn}y=c
\]
\Needspace{24\baselineskip}
woraus folgt, dass $C$ eine Funktion von $c$ ist. Um zu sehen, welche, setzen wir $y=0$; dann wird $C=x$, $c=\operatorname{arc}\sin\operatorname{lemn}x$, also $c=\operatorname{arc}\sin\operatorname{lemn}C$ oder $C=\sin\operatorname{lemn}c$. Ist daher $\sin\operatorname{lemn}p=x$, $\sin\operatorname{lemn}q=y$, so wird
\Needspace{12\baselineskip}
\[
\sin\operatorname{lemn}(p+q)=
\dfrac{x\surd\dfrac{1-yy}{1+yy}+y\surd\dfrac{1-xx}{1+xx}}
{1-xy\surd\dfrac{1-xx}{1+xx}\,\dfrac{1-yy}{1+yy}}
\]

\Needspace{20\baselineskip}
Setzt man nun $p=\tfrac12\varpi$, $q=-a$, so folgt wegen
\Needspace{8\baselineskip}
\[
\sin\operatorname{lemn}(-a)=-\sin\operatorname{lemn}a,\quad
\sin\operatorname{lemn}\dfrac12\varpi=1
\]
\Needspace{10\baselineskip}
\[
\cos\operatorname{lemn}a=\surd\dfrac{1-\sin\operatorname{lemn}a^2}{1+\sin\operatorname{lemn}a^2}
\]
\Needspace{30\baselineskip}
Die vorangehende Form geht aber in die folgende über


\SourceLeaf{424}
\Needspace{14\baselineskip}
\[
\begin{aligned}\sin\operatorname{lemn}(p\pm q)&=\\[7pt]&\quad{}
\dfrac{\sin\operatorname{lemn}p\,\cos\operatorname{lemn}q\pm\sin\operatorname{lemn}q\,\cos\operatorname{lemn}p}
{1\mp\sin\operatorname{lemn}p\,\sin\operatorname{lemn}q\,\cos\operatorname{lemn}p\,\cos\operatorname{lemn}q}\end{aligned}
\]
\Needspace{14\baselineskip}
\[
\begin{aligned}\cos\operatorname{lemn}(p\pm q)&=\\[7pt]&\quad{}
\dfrac{\cos\operatorname{lemn}p\,\cos\operatorname{lemn}q\mp\sin\operatorname{lemn}p\,\sin\operatorname{lemn}q}
{1\pm\sin\operatorname{lemn}p\,\sin\operatorname{lemn}q\,\cos\operatorname{lemn}p\,\cos\operatorname{lemn}q}\end{aligned}
\]

\Needspace{24\baselineskip}
\noindent[Spätere Bemerkung:]

\noindent I.\hfill $\alpha+\beta+\gamma=180^0\;[=\varpi]$\hfill\mbox{}

Setzt man
\Needspace{12\baselineskip}
\[
\begin{array}{ll}
\sin\operatorname{lemn}\alpha=\operatorname{tang}a,&\cos\operatorname{lemn}\alpha=\cos A\\[5pt]
\sin\operatorname{lemn}\beta=\operatorname{tang}b,&\cos\operatorname{lemn}\beta=\cos B\\[5pt]
\sin\operatorname{lemn}\gamma=\operatorname{tang}c,&\cos\operatorname{lemn}\gamma=\cos C
\end{array}
\]
so sind $a,\ b,\ c,\ A,\ B,\ C$ Seiten und Winkel eines sphärischen Dreiecks, in dem
\Needspace{8\baselineskip}
\[
\dfrac{\sin A}{\sin a}=\dfrac{\sin B}{\sin b}=\dfrac{\sin C}{\sin c}=\surd2
\]

\Needspace{24\baselineskip}
\par\Needspace{7\baselineskip}\noindent II.\hfill $\alpha+\beta+\gamma=90^0\;[=\tfrac12\varpi]$\hfill\mbox{}

Setzt man
\Needspace{12\baselineskip}
\[
\begin{array}{ll}
\sin\operatorname{lemn}\alpha=\cos a,&\cos\operatorname{lemn}\alpha=-\operatorname{tang}A\\[5pt]
\sin\operatorname{lemn}\beta=\cos b,&\cos\operatorname{lemn}\beta=-\operatorname{tang}B\\[5pt]
\sin\operatorname{lemn}\gamma=\cos c,&\cos\operatorname{lemn}\gamma=-\operatorname{tang}C
\end{array}
\]
so sind $a,\ b,\ c,\ A,\ B,\ C$ wiederum Seiten und Winkel eines sphärischen Dreiecks, in dem
\Needspace{8\baselineskip}
\[
\dfrac{\sin A}{\sin a}=\dfrac{\sin B}{\sin b}=\dfrac{\sin C}{\sin c}=\surd\dfrac12
\]

\clearpage
\endgroup
